\clearpage\pagestyle{otMain}\setlength{\otBlocEspacement}{2.5mm}\fontsize{11}{15}\selectfont\corrigesChapitre{6}{Découvrir les fractions}
\begin{blocLivre}{cor-6-1}\begin{corrige}{1}{Lire une part}\demarche{Compter les parts du partage puis les parts prises.}
Le disque est partagé en 8 parts égales et 3 sont colorées : la fraction est $\res{\frac38}$.
\label{fig-cor-6-1}\figureGeo[corrige]{frac-1}
\end{corrige}\end{blocLivre}
\begin{blocLivre}{cor-6-2}\begin{corrige}{2}{Colorier une fraction}\demarche{Repérer le nombre de parts égales de la bande.}
La bande compte 5 cases égales : on en colorie $\res2$. La version corrigée montre les deux cases colorées ; tout choix de deux cases convient.
\label{fig-cor-6-2}\figureGeo[corrige]{frac-2}
\end{corrige}\end{blocLivre}
\begin{blocLivre}{cor-6-3}\begin{corrige}{3}{Le bon vocabulaire}\demarche{Le dénominateur nomme la part, le numérateur les compte.}
Le numérateur est $\res7$, le dénominateur $\res{10}$. « Trois quarts » s’écrit $\res{\frac34}$ et « huit sixièmes » $\res{\frac86}$.
\end{corrige}\end{blocLivre}
\begin{blocLivre}{cor-6-4}\begin{corrige}{4}{Des parts vraiment égales ?}\demarche{Contrôler l’égalité des parts avant d’écrire une fraction.}
$\res{\text{Non}}$ : les trois morceaux n’ont pas la même taille. Une fraction ne décrit un partage que si les parts sont égales. Ici le morceau coloré couvre la moitié de la figure.
\label{fig-cor-6-4}\figureGeo[corrige]{frac-4}
\end{corrige}\end{blocLivre}
\begin{blocLivre}{cor-6-5}\begin{corrige}{5}{Au-delà de l’unité}\demarche{Compter les unités complètes dans les parts.}
On compte 5 quarts colorés : $\frac54$. Un disque entier fait 4 quarts : $\res{\frac54=1+\frac14}$.
\label{fig-cor-6-5}\figureGeo[corrige]{frac-5}
\end{corrige}\end{blocLivre}
\begin{blocLivre}{cor-6-6}\begin{corrige}{6}{Entier caché}\demarche{Retirer les paquets qui font des unités entières.}
$\frac74=\frac44+\frac34=\res{1+\frac34}$. Pour les tiers : $\frac93=\res3$, car $9=3\times3$.
\end{corrige}\end{blocLivre}
\begin{blocLivre}{cor-6-7}\begin{corrige}{7}{Recomposer}\demarche{Convertir les entiers dans la part demandée.}
Deux unités font $\frac{10}5$. Avec les trois cinquièmes restants : $2+\frac35=\res{\frac{13}5}$.
\end{corrige}\end{blocLivre}
\begin{blocLivre}{cor-6-8}\begin{corrige}{8}{Encadrer une fraction}\demarche{Chercher les multiples du dénominateur voisins du numérateur.}
$\frac84=2$ et $\frac{12}4=3$ : $\res{2<\frac{11}4<3}$. De même $\frac{15}5=3$ et $\frac{20}5=4$ : $\res{3<\frac{17}5<4}$.
\end{corrige}\end{blocLivre}
\begin{blocLivre}{cor-6-9}\begin{corrige}{9}{Sept galettes pour quatre}\demarche{Couper chaque galette selon le nombre d’enfants.}
Chaque enfant reçoit $\res{\frac74}$ de galette : une galette entière et trois quarts, soit $\res{1+\frac34}$. Contrôle : $4\times\frac74=7$ galettes.
\end{corrige}\end{blocLivre}
\begin{blocLivre}{cor-6-10}\begin{corrige}{10}{Égales à un entier}\demarche{Une fraction vaut 1 quand numérateur et dénominateur sont égaux.}
Il faut autant de parts que le partage : $\res{\frac66=1}$ et $\res{\frac{12}{12}=1}$. Deux unités font $2\times4$ quarts, soit $\res{\frac84}$.
\end{corrige}\end{blocLivre}
\begin{blocLivre}{cor-6-11}\begin{corrige}{11}{Le gâteau de la fête}\demarche{Compter les parts sur le même partage en 12.}
On a mangé $5+4=9$ parts : $\res{\frac9{12}}$ du gâteau. Il reste $12-9=3$ parts, soit $\res{\frac3{12}}$. Contrôle : $\frac9{12}+\frac3{12}=\frac{12}{12}=1$, le gâteau entier.
\label{fig-cor-6-11}\figureGeo[corrige]{frac-11}
\end{corrige}\end{blocLivre}
\begin{blocLivre}{cor-6-12}\begin{corrige}{12}{La frise de la classe}\demarche{Prendre une part du partage en quatre.}
Le quart de 20 : $20:4=\res5$ carreaux bleus. Il reste 15 carreaux, soit $\res{\frac{15}{20}}$ de la frise, les trois quarts.
\label{fig-cor-6-12}\figureGeo[corrige]{frac-12}
\end{corrige}\end{blocLivre}
\begin{blocLivre}{cor-6-13}\begin{corrige}{13}{Des tartes en sixièmes}\demarche{Comparer le numérateur au dénominateur.}
Sept sixièmes : $\frac76$. Comme $7>6$, cette fraction $\res{\text{dépasse }1}$ : une tarte entière et un sixième, $\res{\frac76=1+\frac16}$.
\end{corrige}\end{blocLivre}
\begin{blocLivre}{cor-6-14}\begin{corrige}{14}{Dépasser 1 sans calculer}\demarche{Comparer numérateur et dénominateur.}
$\frac98$ et $\frac{20}{19}$ sont $\res{\text{supérieures à }1}$ : le numérateur dépasse le dénominateur. $\frac56<1$ et $\frac{12}{12}=1$.
\end{corrige}\end{blocLivre}
\begin{blocLivre}{cor-6-15}\begin{corrige}{15}{La moitié de qui ?}\demarche{Toujours préciser le tout dont on prend une fraction.}
$\res{\text{Léna a raison}}$ : une fraction désigne une part d’un tout précis. La moitié d’une petite pizza peut être moins grosse que le quart d’une pizza géante. Sur un même tout, en revanche, $\frac12>\frac14$.
\end{corrige}\end{blocLivre}
